\section{Preliminaries and Problem Model}
\label{sec:preliminaries}
\subsection{Notation and Public Parameters}
\label{sec:notation-params}
For a positive integer \(d\), write \([d]=\{1,\ldots,d\}\).
The security parameter is \(\kappa\); PPT means probabilistic polynomial time.
For an integer \(z\) and modulus \(d\ge2\), \([z]_d\) is its canonical
representative in \(\{0,\ldots,d-1\}\). This notation extends coefficient-wise.
For a monic polynomial \(f\in\mathbb Z[X]\) of degree \(n\), write
\[
 R=\mathbb Z[X]/(f),\qquad R_q=R/qR,\qquad R_t=R/tR.
\]
The construction uses \(f=X^n+1\), with \(n\) a power of two,
\(t=2^w\), \(q>t\), and \(\gcd(t,q)=1\).
The privacy reduction does not require \(q\) to be prime.
All parameter descriptions have polynomial bit length in \(\kappa\);
ring degrees and sample counts in the reductions are polynomially bounded.
Every ring element has a unique representative of degree less than \(n\).
For \(P\in R_t\), \(\widehat P\in R\) denotes the lift whose coefficients
lie in \([0,t)\). This lift, followed by reduction modulo \(q\), is an
encoding into \(R_q\), not a ring homomorphism from \(R_t\).
For \(z\in R_q\), \(\operatorname{center}_q(z)\in R\) is obtained by choosing
each coefficient in \([-q/2,q/2)\). The infinity norm in \(R\) is the maximum
absolute value of its integer coefficients.

\begin{table}[t]
\centering
\caption{Public parameters and encoding units.}
\label{tab:notation}
\small
\begin{tabular}{@{}lp{8.9cm}@{}}
\toprule
Symbol & Meaning \\
\midrule
\(N,n\) & Number of database records; coefficients per polynomial (\(n\) a power of two). \\
\(\ell,\rho_0\) & Original record length and segment width, in bits; both at least one. \\
\(B,w,t\) & Segment bound \(B=2^{\rho_0}\); selected backend coefficient width \(w\); \(t=2^w\). \\
\(J,L\) & Segments per record \(J=\lceil\ell/\rho_0\rceil\); polynomial blocks per record \(L=\lceil J/n\rceil\). \\
\(\alpha,K\) & Number of distinct targets, \(1\le\alpha\le N\); derived radix weights \(K=(1,B,\ldots,B^{\alpha-1})\). \\
\(\ell_{\rm enc}\) & Full-block encoded bit length \(Lnw\); distinct from \(\ell\) and \(Jw\). \\
\(\chi_r,r\) & Coefficient-wise rounded Gaussian family and its public continuous standard deviation \(r>0\); \(\chi_s=\chi_e=\chi_r\). \\
\(p_{\rm unit}\) & Probability that a uniform element of \(R_q\) is invertible. \\
\bottomrule
\end{tabular}
\end{table}

The public parameters \(pp\) specify \(\kappa,N,f,n,q,t,w,\ell,\rho_0,\alpha,K\),
the public width \(r>0\), the family \(\chi=\chi_r\) defined below, and the
encoding layout; \(B,K,J,L\) are derived, with \(k_\nu=B^{\nu-1}\), \(\nu\in[\alpha]\). We retain
\(\chi_s=\chi_e=\chi_r\) as aliases in the construction and its separate
correctness statement. All public parameters are fixed independently
of the secret target tuple. In particular, \(pp\)
is not a public encryption key. The exact radix capacity condition is
\(t\ge B^\alpha\), equivalently \(w\ge\alpha\rho_0\), as in
Lemma~\ref{lemma:packing-condition}. The weights are canonical elements of
\(\mathbb Z_t\), not necessarily units; no weight coprimality or inverse
is required. The condition \(\gcd(t,q)=1\) remains necessary for the
security reduction. 

\subsection{Retrieval Task and Interfaces}
\label{sec:mq-model}
A legal database is an ordered list \(M=(m_1,\ldots,m_N)\in
(\{0,1\}^{\ell})^N\). A legal request is an ordered tuple
\(I=(i_1,\ldots,i_\alpha)\in[N]^\alpha\) of pairwise-distinct indices;
record contents at distinct indices may be equal. Write
\(M[I]=(m_{i_1},\ldots,m_{i_\alpha})\).
The deterministic record map \(\mathsf{Encode}:\{0,1\}^{\ell}\to R_t^L\)
is specified in Section~\ref{sec:encoding}. We use
\(\mathsf{DBEncode}(M)=(\mathsf{Encode}(m_i))_{i\in[N]}\)
only as its row-wise abbreviation, not as an additional protocol.

\begin{table}[t]
\centering\small
\caption{Protocol interfaces and ownership. Public parameters determine all
dimensions. Query generation is randomized; reply generation and extraction
are deterministic for their inputs.}
\label{tab:interfaces}
\begin{tabularx}{\linewidth}{@{}llX@{}}
\toprule Algorithm & Party & Input and output \\
\midrule
\(\mathsf{KeyGen}\) & Client & \(pp\mapsto sk=s\); one fresh secret per query. \\
\(\mathsf{QueryGen}\) & Client & \((pp,sk,I)\mapsto(Q,\mathsf{st})\), where
\(Q\in(R_q^2)^N\) and private \(\mathsf{st}=I\). \\
\(\mathsf{ReplyGen}\) & Server & \((pp,P,Q)\mapsto\mathsf{Resp}\in(R_q^2)^L\),
where \(P=\mathsf{DBEncode}(M)\in R_t^{N\times L}\). \\
\(\mathsf{ReplyExt}\) & Client & \((pp,sk,\mathsf{Resp},\mathsf{st})\mapsto
\widehat M\in(\{0,1\}^{\ell})^\alpha\), in the requested order. \\
\bottomrule
\end{tabularx}
\end{table}
Only \(Q\) leaves the client; \(sk\) and \(\mathsf{st}\) remain private.
A single query is a vector of \(N\) ciphertexts, and the response contains
\(L\) ciphertexts. Public parameters, including the encoding and weights,
are fixed independently of the target tuple.

\subsection{Correctness Goal}
For every admissible \(pp\), legal \(M\), and legal \(I\), consider the
honest execution
\[
 sk\leftarrow\mathsf{KeyGen}(pp),\quad
 (Q,\mathsf{st})\leftarrow\mathsf{QueryGen}(pp,sk,I),
\]
\[
 \mathsf{Resp}=\mathsf{ReplyGen}(pp,\mathsf{DBEncode}(M),Q),\quad
 \widehat M=\mathsf{ReplyExt}(pp,sk,\mathsf{Resp},\mathsf{st}).
\]
Define \(\mathrm{Corr}=[\widehat M=M[I]]\). A correctness target
\(\Pr[\neg\mathrm{Corr}]\le\delta_{\rm corr}(pp)\) quantifies over all
such databases and requests, with probability over key generation and
all encryption randomness. This target is a specification, not an
unconditional negligible-failure claim. Section~\ref{sec:corr} proves
ordered recovery under an integer no-wrap condition and derives a bound
for the specified error law.

\subsection{Single-challenge Index Privacy}
\label{sec:mq-privacy}
\begin{definition}[Index privacy for packed multi-record retrieval]
\label{def:mq-cpir-privacy}
Fix admissible public parameters \(pp\) independently of the targets. The database and target tuples must be legal as defined above; both
challenge tuples use the same public layout and weight tuple. For a two-stage PPT adversary \(\mathcal A=(\mathcal A_1,\mathcal A_2)\),
the experiment is:
\begin{enumerate}
\item The challenger samples \(sk\leftarrow\mathsf{KeyGen}(pp)\).
\item \(\mathcal A_1(pp)\) outputs \((M,I_0,I_1,\sigma)\), where \(\sigma\)
is its polynomial-length private state. Both tuples and the database must be
legal. The tuples may coincide, covering degenerate instances such as
\(N=\alpha=1\). On illegal input, the experiment terminates with an independent
fair success bit.
\item The challenger samples \(b\leftarrow\{0,1\}\), computes
\[
 (Q,\mathsf{st})\leftarrow\mathsf{QueryGen}(pp,sk,I_b),\qquad
 P=\mathsf{DBEncode}(M),
\]
\[
 \mathsf{Resp}=\mathsf{ReplyGen}(pp,P,Q),
\]
and gives \((Q,\mathsf{Resp})\) to the adversary.
\item \(b'\leftarrow\mathcal A_2(pp,M,I_0,I_1,\sigma,Q,\mathsf{Resp})\).
The experiment succeeds iff \(b'=b\).
\end{enumerate}
Neither \(sk\) nor the client state \(\mathsf{st}\) is disclosed. Define
\[
 \Adv^{\rm priv}_{\Pi,\mathcal A}(\kappa)
   =\left|\Pr[\text{experiment succeeds}]-\tfrac12\right|.
\]
Single-challenge index privacy means this advantage is negligible for every
PPT adversary. The probability includes key generation, all encryptions,
the challenge bit, and adversary randomness.
\end{definition}
The public parameters, database, and adversary state are part of the server's
view. Its deterministic intermediate computations are functions of these
values and the query. No decryption outcome or client extraction state is
included in this game.

\subsection{Specified Error Law}
We adopt the coefficient-wise Gaussian family explicitly defined by
Brakerski and Vaikuntanathan~\cite[Sec.~1.3]{homomorphic}:
\[
 Y_j\leftarrow\mathcal N(0,r^2),\qquad
 Z_j=\operatorname{round}(Y_j),\qquad
 \chi_r=\mathcal L\!\left(\sum_{j=0}^{n-1}Z_jX^j\right),
\]
where the \(Y_j\) are independent and rounding is to the nearest integer.
Equivalently, one coefficient has probability mass
\[
 p_r(z)=\int_{z-1/2}^{z+1/2}
       \frac{\exp(-y^2/(2r^2))}{\sqrt{2\pi}\,r}\,dy,\qquad z\in\mathbb Z.
\]
Half-integer ties have probability zero under the ideal continuous law.
Both the secret and every fresh encryption error use this same
product distribution, independently across coefficients and draws.
Here \(r\) is the standard deviation before rounding; its numerical value
remains to be selected jointly with the other parameters.


This rounded law differs from the integer Gaussian with mass proportional
to \(\exp(-\pi z^2/\eta^2)\) and is not identified here with a
canonical-embedding Gaussian. Efficient machine realization and its
approximation to this ideal law require separate justification.

\subsection{Linear Encryption Interface}
\label{sec:prelim}
We use the symmetric-key algebraic interface of the Ring-LWE construction
of Brakerski and Vaikuntanathan~\cite{homomorphic}, in the linear evaluation
setting of XPIR~\cite{melchor2016xpir}. All ciphertext arithmetic below is in
\(R_q\); all message lifts use canonical coefficients.
\begin{align*}
 \mathsf{KeyGen}(pp)&:\ s\leftarrow\chi,\quad sk=s,\\
 \mathsf{Enc}_s(v)&:\ a\leftarrow R_q\text{ uniformly},\quad e\leftarrow\chi,
   \quad C=(as+te+\widehat v,-a),\\
 \mathsf{Dec}_s(C)&:=
   [\operatorname{center}_q(c_0+c_1s)]_t,\\
 \mathsf{Add}(C,C')&:=(c_0+c'_0,c_1+c'_1),\\
 \Absorb(P,C)&:=(\widehat P c_0,\widehat P c_1).
\end{align*}
Each encryption samples fresh independent \(a,e\), including encryptions of
zero. The client alone possesses \(sk\), encrypts queries, and decrypts
replies. The server needs only \(pp\), the encoded database, and ciphertexts
to perform \(\mathsf{Add}\) and \(\Absorb\). No ciphertext--ciphertext
multiplication or public key is required.

XPIR's mathematical representation \((a,b)\), with \(b=as+te+v\), corresponds
to our \(C=(b,-a)\), so \(c_0+c_1s=b-as\). This correspondence makes no claim
of identical binary serialization. The final mod-\(t\) step returns canonical
coefficients, whereas the preceding mod-\(q\) step uses centered coefficients.
Correctness requires control of the integer value before reduction modulo
\(q\), as formalized in Theorem~\ref{thm:correctness}.

\subsection{Uniform-secret Decisional PLWE}
\label{sec:rlwe-plwe}
\begin{definition}[Uniform-secret decisional PLWE]
\label{def:uplwe}
Fix a publicly specified parameter family \((f,q,\chi)\), where \(\chi\)
is efficiently sampleable over \(R\). For a polynomially bounded
\(\tau=\tau(\kappa)\), consider
\[
 \mathcal U_{\rm real}^{(\tau)}
   =((a_i,a_i s+e_i))_{i=1}^{\tau},\qquad
 \mathcal U_{\rm unif}^{(\tau)}
   =((a_i,u_i))_{i=1}^{\tau}
\]
in \((R_q^2)^\tau\). One secret \(s\leftarrow U(R_q)\) is shared by all
samples. All \(a_i,u_i\leftarrow U(R_q)\) and \(e_i\leftarrow\chi\)
are mutually independent and independent of \(s\); errors are reduced
modulo \(q\) when used in \(R_q\). For a PPT distinguisher \(\mathcal B\), define
\[
 \Adv^{\rm uPLWE}_{\mathcal B}(\kappa;\tau)=
 \left|\Pr[\mathcal B(pp,\mathcal U_{\rm real}^{(\tau)})=1]
       -\Pr[\mathcal B(pp,\mathcal U_{\rm unif}^{(\tau)})=1]\right|.
\]
The \(\tau\)-sample assumption for this family is that this advantage is
negligible in \(\kappa\) for every PPT \(\mathcal B\).
\end{definition}
This is a distribution-specific decisional assumption, following the
uniform-secret PLWE formulation discussed in~\cite[Sec.~2]{homomorphic}.
Hardness is assumed only for the corresponding instance, not for every
efficiently sampleable \(\chi\). We make no worst-case ideal-lattice
reduction claim and certify no concrete security level.
